\chapter{Hamiltonian Evolution and Time-Sliced Path Integrals}
\label{ch:hamiltonian-evolution-time-sliced-path-integrals}
\ConventionsRealTime


\section{Classical Lagrangian and Hamiltonian Data}

\begin{definition}[Regular classical Lagrangian datum]
\label{def:regular-classical-lagrangian-datum}
Let \(Q\) be a smooth \(d\)-dimensional configuration manifold.  Local
coordinates on \(Q\) are denoted by \(q^a\), \(a=1,\ldots,d\).  A tangent
vector is written \((q,\dot q)\in TQ\), and a cotangent vector is written
\((q,p)\in T^*Q\), with \(p=p_a\,\dd q^a\).  A classical Lagrangian is a
function
\[
  L:TQ\to\mathbb R.
\]
For a path \(q:[t_i,t_f]\to Q\), the action is
\[
  S[q]=\int_{t_i}^{t_f}L(q(t),\dot q(t))\,\dd t .
\]
The variational principle with fixed endpoints gives the Euler--Lagrange
equations
\[
  \frac{\dd}{\dd t}\frac{\partial L}{\partial \dot q^a}
  =
  \frac{\partial L}{\partial q^a}.
\]
Thus a Lagrangian description uses the pair \((q,\dot q)\) as local Cauchy
data for the classical second-order equation.

Assume now that the fiber derivative of \(L\),
\[
  (q,\dot q)\longmapsto
  \left(q,p_a=\frac{\partial L}{\partial\dot q^a}\right),
\]
is invertible on the region under consideration.  The Hamiltonian associated
to this regular Lagrangian is the Legendre transform
\[
  H(q,p)=p_a\dot q^a-L(q,\dot q)
  \quad
  \text{with }\dot q=\dot q(q,p).
\]
The canonical Poisson bracket on \(T^*Q\) is
\[
  \{f,g\}_{\rm P}
  =
  \frac{\partial f}{\partial q^a}\frac{\partial g}{\partial p_a}
  -
  \frac{\partial f}{\partial p_a}\frac{\partial g}{\partial q^a},
\]
where repeated indices are summed.  Hamiltonian time evolution is
\[
  \dot f=\{f,H\}_{\rm P},
\]
and, in particular,
\[
  \dot q^a=\frac{\partial H}{\partial p_a},
  \qquad
  \dot p_a=-\frac{\partial H}{\partial q^a}.
\]
The Lagrangian and Hamiltonian descriptions are equivalent only after the
regularity hypothesis above has been imposed; without it, constraints must be
included as part of the classical data.
\end{definition}

\section{Canonical Quantum-Mechanical Data}

\begin{definition}[Schrodinger representation datum]
\label{def:schrodinger-representation-datum}
Let \(d\) be the number of degrees of freedom and let
\(\mathcal Q=\mathbb R^d\) be the configuration space. The Hilbert space is
\[
  \Hilb_{\mathcal Q}=L^2(\mathcal Q,\dd^dq).
\]
The coordinate operators \(\widehat q^a\), \(a=1,\ldots,d\), act by
multiplication. The momentum operators are
\[
  \widehat p_a=-\ii\hbar\,\frac{\partial}{\partial q^a}
\]
first on the common invariant test domain
\(\mathcal S(\mathbb R^d)\subset\Hilb_{\mathcal Q}\), and satisfy
\[
  [\widehat q^a,\widehat p_b]=\ii\hbar\,\delta^a_b.
\]
Closures or self-adjoint extensions are specified when a Hamiltonian is
chosen.  The Schwartz-domain convention is enough for the distributional kernel
manipulations below; the operator limit is then stated at the Hilbert-space
level.

Let \(\widehat H\) be a self-adjoint Hamiltonian on \(\Hilb_{\mathcal Q}\).
This is an analytic hypothesis on the chosen Hilbert-space realization of the
formal classical Hamiltonian.  By Stone's theorem, self-adjointness of
\(\widehat H\) is exactly the condition that the following formula defines a
strongly continuous unitary one-parameter group with generator
\(-\ii\widehat H/\hbar\).  The time-evolution operator over
time \(T\in\mathbb R\) is
\[
  U(T)=\exp\left(-\frac{\ii}{\hbar}T\widehat H\right).
\]
In the Schr\"odinger picture, a vector evolves as
\[
  \ket\psi\longmapsto U(t)\ket\psi .
\]
In the Heisenberg picture, an operator \(\widehat O\) evolves as
\[
  \widehat O(t)=U(t)^{-1}\widehat O\,U(t),
  \qquad
  \frac{\dd}{\dd t}\widehat O(t)
  =
  \frac{\ii}{\hbar}[\widehat H,\widehat O(t)]
\]
on vectors for which the commutator expression is defined.  The path integral
constructed below gives a regulated representation of matrix elements of this
Hamiltonian time evolution.
\end{definition}

\begin{definition}[Position-momentum rigging and kernel notation]
\label{def:position-momentum-rigging-kernel-notation}
The position and momentum generalized eigenstates are denoted by
\(\ket{q}\) and \(\ket{p}\). They are elements of
\(\mathcal S'(\mathbb R^d)\) in the rigged Hilbert space
\[
  \mathcal S(\mathbb R^d)
  \subset
  L^2(\mathbb R^d,\dd^dq)
  \subset
  \mathcal S'(\mathbb R^d),
\]
with the \(\hbar\)-dependent Fourier convention
\[
  \widetilde\psi(p)
  =
  \int_{\mathbb R^d}\dd^dq\,
  \ee^{-\frac{\ii}{\hbar}p_aq^a}\psi(q).
\]
They are normalized by
\[
  \braket{q}{p}=\ee^{\frac{\ii}{\hbar}p_aq^a},
\]
and used through the resolutions
\[
  \int_{\mathbb R^d}\dd^dq\,\ket{q}\bra{q}=1,
  \qquad
  \int_{\mathbb R^d}\frac{\dd^dp}{(2\pi\hbar)^d}\,\ket{p}\bra{p}=1.
\]
These are weak identities: after inserting test vectors on the left and right
they become the ordinary Fourier inversion formula.
The transition amplitude is the distributional kernel
\[
  K(q_f,T;q_i,0)
  =
  \bra{q_f}U(T)\ket{q_i}.
\]
The continuum Lagrangian path-integral notation for this same kernel is
\[
  \bra{q_f}U(T)\ket{q_i}
  =
  \int_{q(0)=q_i}^{q(T)=q_f} [Dq]\,
  \exp\left(\frac{\ii}{\hbar}S[q]\right).
\]
This displayed equation is a compact name for the regulated construction
below.  The bracketed symbols \([Dq]\) and, in phase-space formulas, \([Dp]\)
are determined by the regulator, limiting procedure, operator symbol, and
boundary conditions; at fixed regulator the object is a finite-dimensional
oscillatory integral.  Unbracketed products such as \(\dd^dq_n\) and
\(\dd^dp_n\) are reserved for the finite-dimensional measures that define the
regulated expression.
\end{definition}

\section{Finite Time Slicing}

\paragraph{Finite phase-space time-sliced construction.}
\label{par:finite-phase-space-time-sliced-kernel}
Fix the data of Definition~\ref{def:position-momentum-rigging-kernel-notation}
and an operator symbol \(h(q,p)\) for \(\widehat H\) in a specified ordering
convention.  The finite-\(N\) phase-space expression \(K_N\) is the regulated
oscillatory integral obtained by inserting position and momentum resolutions
in the product formula for \(U(T)\).  The theorem-level analytic input enters
later, through the operator convergence theorem for the limiting procedure.
We record the construction explicitly.
Choose an integer \(N\ge1\), set \(\epsilon=T/N\), and define
\[
  t_n=n\epsilon,\qquad n=0,\ldots,N,
\]
with \(q_0=q_i\) and \(q_N=q_f\).  Inserting \(N-1\) position resolutions and
\[
  U(T)=\left[\exp\left(-\frac{\ii}{\hbar}\epsilon\widehat H\right)\right]^N
\]
gives first the position-space finite product
\[
  K(q_f,T;q_i,0)
  =
  \int\prod_{n=1}^{N-1}\dd^dq_n\,
  \prod_{n=0}^{N-1}
  \bra{q_{n+1}}
  \exp\left(-\frac{\ii}{\hbar}\epsilon\widehat H\right)
  \ket{q_n}.
\]
Inserting \(N\) momentum resolutions gives
\[
\begin{aligned}
  K(q_f,T;q_i,0)
  &=
  \int\prod_{n=1}^{N-1}\dd^dq_n
  \prod_{n=0}^{N-1}
  \frac{\dd^dp_n}{(2\pi\hbar)^d}
  \\
  &\quad\times
  \prod_{n=0}^{N-1}
  \bra{q_{n+1}}p_n\rangle
  \bra{p_n}\exp\left(-\frac{\ii}{\hbar}\epsilon\widehat H\right)\ket{q_n}.
\end{aligned}
\]

\begin{center}
\begin{tikzpicture}[>=Stealth]
  \draw[line width=0.6pt] (0,0) -- (9.2,0);
  \foreach \x/\lab in {
    0/{q_0=q_i},
    1.25/{q_1},
    2.5/{q_2},
    3.75/{q_3},
    6.75/{q_{N-1}},
    8.75/{q_N=q_f}
  }{
    \draw (\x,0.11) -- (\x,-0.11);
    \node[above] at (\x,0.11) {$\lab$};
  }
  \node[below] at (0,-0.12) {$t_0=0$};
  \node[below] at (8.75,-0.12) {$t_N=T$};
  \node at (5.2,0.19) {$\cdots$};
  \draw[decorate,decoration={brace,amplitude=4pt,mirror}]
    (0,-0.5) -- node[below=3pt] {$\epsilon=T/N$} (1.25,-0.5);
\end{tikzpicture}
\end{center}

Suppose a symbol \(h(q,p)\) for \(\widehat H\) has been fixed by an operator
ordering convention, for example by
\[
  \bra{p}\widehat H\ket{q}=h(q,p)\braket{p}{q}
\]
on the chosen domain.  Replacing each short-time matrix element by the first
terms of its symbol expansion gives
\[
  \bra{q_{n+1}}p_n\rangle
  \bra{p_n}\ee^{-\ii\epsilon\widehat H/\hbar}\ket{q_n}
  =
  \exp\left[
    \frac{\ii}{\hbar}
    \left(
      p_{n,a}(q_{n+1}^a-q_n^a)
      -\epsilon h(q_n,p_n)
      +O(\epsilon^2)+O(\hbar\epsilon)
    \right)
  \right]
\]
for this convention, with the error understood in the regulated
operator-symbol expansion rather than as a pointwise statement about an
undefined continuum integral.  The associated finite-\(N\) phase-space
approximant is
\[
  K_N(q_f,T;q_i,0)
  =
  \int
  \prod_{n=1}^{N-1}\dd^dq_n
  \prod_{n=0}^{N-1}\frac{\dd^dp_n}{(2\pi\hbar)^d}
  \exp\left(\frac{\ii}{\hbar}S_N[p,q]\right),
\]
where the discrete phase-space action is
\[
  S_N[p,q]
  =
  \sum_{n=0}^{N-1}
  \left[
    p_{n,a}(q_{n+1}^a-q_n^a)
    -\epsilon\,h(q_n,p_n)
  \right].
\]
The continuum notation
\[
  \int [Dq]\,[Dp]\,
  \exp\left[
    \frac{\ii}{\hbar}\int_0^T
    \bigl(p_a\dot q^a-h(q,p)\bigr)\dd t
  \right]
\]
denotes this regulated construction together with the specified limiting
procedure. In the present chapter, the finite-\(N\) expression is the defined
object.

\section{Trotter Limits, Wiener Measure, and Kernel Convergence}

\begin{theorem}[Trotter product formula in the form used here]
\label{qthm:trotter-product-formula-qm}
The passage from the finite product to a continuum path integral is a
convergence statement about operators before it is a statement about kernels.
Let \(\widehat A\) and \(\widehat B\) be self-adjoint operators on
\(\Hilb_{\mathcal Q}\).  For the real-time statement, assume that
\(\operatorname{Dom}\widehat A\cap\operatorname{Dom}\widehat B\) is dense and
that the algebraic sum
\[
  \widehat A+\widehat B
\]
on this domain is essentially self-adjoint, with self-adjoint closure
\(\widehat H\).  Then
\[
  \exp\left(-\frac{\ii T}{\hbar}\widehat H\right)
  =
  \operatorname*{s-lim}_{N\to\infty}
  \left[
    \exp\left(-\frac{\ii\epsilon}{\hbar}\widehat A\right)
    \exp\left(-\frac{\ii\epsilon}{\hbar}\widehat B\right)
  \right]^N,
  \qquad
  \epsilon=\frac{T}{N},
\]
where the limit is in the strong operator topology on \(\Hilb_{\mathcal Q}\).
For the Euclidean statement, assume instead that \(\widehat A\) and
\(\widehat B\) are lower-semibounded and that their closed quadratic forms
\(\mathfrak a\) and \(\mathfrak b\) have dense common form domain
\[
  \mathcal Q_{\mathfrak h}
  =
  \operatorname{Dom}\mathfrak a\cap\operatorname{Dom}\mathfrak b ,
\]
with \(\mathfrak h=\mathfrak a+\mathfrak b\) closed and lower-semibounded on
\(\mathcal Q_{\mathfrak h}\).  If \(\widehat H=\widehat A\dotplus\widehat B\)
is the self-adjoint form sum associated to \(\mathfrak h\), then
\[
  \exp\left(-\frac{\tau}{\hbar}\widehat H\right)
  =
  \operatorname*{s-lim}_{N\to\infty}
  \left[
    \exp\left(-\frac{\delta}{\hbar}\widehat A\right)
    \exp\left(-\frac{\delta}{\hbar}\widehat B\right)
  \right]^N,
  \qquad
  \delta=\frac{\tau}{N},
\]
again in the strong operator topology.
\end{theorem}

\begin{proof}
The product formula is a statement in \(\Hilb_{\mathcal Q}\), not a pointwise
identity between formal path-integral densities.  In the bounded case the
mechanism can be seen without domain issues.  If \(A\) and \(B\) are bounded
self-adjoint operators and \(t\) is fixed, Taylor expansion in operator norm
gives, uniformly for \(0\le\delta\le t\),
\[
  \ee^{-\delta A}\ee^{-\delta B}
  =
  1-\delta(A+B)+O(\delta^2),
  \qquad
  \ee^{-\delta(A+B)}
  =
  1-\delta(A+B)+O(\delta^2).
\]
Hence
\[
  \bigl\|\ee^{-\delta A}\ee^{-\delta B}
  -\ee^{-\delta(A+B)}\bigr\|\le C_t\delta^2 .
\]
The telescoping identity
\[
  X^N-Y^N
  =
  \sum_{j=0}^{N-1}X^{N-1-j}(X-Y)Y^j
\]
with \(X=\ee^{-tA/N}\ee^{-tB/N}\) and
\(Y=\ee^{-t(A+B)/N}\) then gives
\[
  \|X^N-Y^N\|
  \le
  N\,\|X-Y\|\,\max_{j,k\le N}\|X^j\|\,\|Y^k\|
  =
  O(N^{-1}).
\]
Thus the noncommutativity of \(A\) and \(B\) first appears in terms of order
\(\delta^2\), and the accumulated error over \(N=t/\delta\) time slices
vanishes.

For Schr\"odinger Hamiltonians \(T+V\) the same calculation cannot be read as
a proof, because \(T\) and \(V\) are unbounded and may not have a useful
operator-domain intersection.  The semibounded Euclidean statement is instead
form-theoretic.  After adding constants to make \(A\) and \(B\) nonnegative
if necessary--which only multiplies the product formula by the corresponding
scalar exponential factors--one starts from closed semibounded quadratic forms
\[
  \mathfrak a[\psi]=\langle A^{1/2}\psi,A^{1/2}\psi\rangle,\qquad
  \mathfrak b[\psi]=\langle B^{1/2}\psi,B^{1/2}\psi\rangle ,
\]
with common form domain
\(\operatorname{Dom}\mathfrak a\cap\operatorname{Dom}\mathfrak b\).  If
\(\mathfrak h=\mathfrak a+\mathfrak b\) is closed and semibounded, the
representation theorem for closed forms defines the self-adjoint form sum
\(H=A\dotplus B\) by
\[
  \mathfrak h(\varphi,\psi)=\langle\varphi,H\psi\rangle,
  \qquad
  \psi\in\operatorname{Dom} H,\quad
  \varphi\in\operatorname{Dom}\mathfrak h .
\]
The Trotter--Kato form theorem proves that the alternating bounded
contractions
\[
  \bigl(\ee^{-tA/N}\ee^{-tB/N}\bigr)^N
\]
converge strongly to \(\ee^{-tH}\).  Its proof replaces the norm Taylor
argument by resolvent and form convergence: the resolvent of \(H\) is
characterized by the variational equation
\[
  \mathfrak h(\varphi,u)+\lambda\langle\varphi,u\rangle
  =
  \langle\varphi,f\rangle,
  \qquad
  \varphi\in\operatorname{Dom}\mathfrak h,
\]
and the product approximants have resolvents whose variational problems
converge to this one on the common form domain.  Strong resolvent convergence
then gives strong semigroup convergence by the spectral theorem.

For the real-time statement, define
\[
  F(s)=\exp\left(-\frac{\ii s}{\hbar}A\right)
       \exp\left(-\frac{\ii s}{\hbar}B\right).
\]
On the dense algebraic core
\(\operatorname{Dom}A\cap\operatorname{Dom}B\),
\[
  \lim_{s\to0}\frac{F(s)\psi-\psi}{s}
  =
  -\frac{\ii}{\hbar}(A+B)\psi .
\]
The factors are unitary, hence \(\|F(s)\|\le1\).  Chernoff's product theorem
then identifies the strong limit of \(F(T/N)^N\) with the unitary group whose
generator is the closure of \(-\ii(A+B)/\hbar\), namely
\(\exp(-\ii T H/\hbar)\).  This real-time conclusion is therefore a theorem
about self-adjoint generators; it should not be inferred by simply
substituting \(t\mapsto\ii t\) into an uncontrolled kernel expression.
\end{proof}

\begin{definition}[Kato Schrodinger Euclidean datum]
\label{def:kato-schrodinger-euclidean-datum}
For Schr\"odinger operators the words ``suppose the hypotheses hold'' can be
replaced by concrete analytic criteria.  Work on
\(L^2(\mathbb R^d,\dd^dq)\), let
\[
  \widehat p_a=-\ii\hbar\partial_a,
  \qquad
  T=\frac{1}{2m}\widehat p^2=-\frac{\hbar^2}{2m}\Delta ,
\]
and let \(V(\widehat q)\) be multiplication by a real measurable function
\(V\).  For a measurable function \(W\), define the Kato class
\(\mathcal K_d\) by
\[
  \lim_{t\downarrow0}
  \sup_{x\in\mathbb R^d}
  \int_0^t
  \bigl(\ee^{s\Delta}|W|\bigr)(x)\,\dd s
  =
  0,
\]
where \(\ee^{s\Delta}\) is the heat semigroup on \(\mathbb R^d\).  Define the
local Kato class by requiring \(\chi_K W\in\mathcal K_d\) for every compact
\(K\subset\mathbb R^d\).  A standard sufficient Euclidean hypothesis is the
Kato decomposition
\[
  V=V_+-V_-,
  \qquad
  V_\pm\ge0,
  \qquad
  V_+\in\mathcal K_{d,\operatorname{loc}},
  \qquad
  V_-\in\mathcal K_d,
\]
together with a lower bound when the literal multiplication factors
\(\ee^{-\delta V/\hbar}\) are used as bounded operators.  Under these
hypotheses the negative part is infinitesimally form-bounded relative to the
free kinetic form: for every \(a>0\) there is \(b_a<\infty\) such that
\[
  \int_{\mathbb R^d} V_-(q)|\psi(q)|^2\,\dd^dq
  \le
  a\int_{\mathbb R^d}|\nabla\psi|^2\,\dd^dq
  +
  b_a\int_{\mathbb R^d}|\psi(q)|^2\,\dd^dq .
\]
Therefore the KLMN theorem applies.  The quadratic form
\[
  Q[\psi]
  =
  \frac{\hbar^2}{2m}\int_{\mathbb R^d}|\nabla\psi|^2\,\dd^dq
  +
  \int_{\mathbb R^d}V(q)|\psi(q)|^2\,\dd^dq
\]
on \(H^1(\mathbb R^d)\cap L^2(V_+(q)\dd^dq)\) is closed and semibounded, and it
defines the Friedrichs Hamiltonian \(H=T+V\).  The Euclidean semigroup has the
Trotter--Kato product formula, and its kernel is represented by the
Feynman--Kac construction whenever the stated Kato hypotheses hold.
\end{definition}

\begin{theorem}[Wiener measure and the Feynman--Kac formula]
\label{thm:wiener-feynman-kac-qm}
Let \(V:\mathbb R^d\to\mathbb R\) be real, locally Kato, and bounded from
below.  Let \(H\) be the Friedrichs Hamiltonian associated to
\[
  Q[\psi]
  =
  \frac{\hbar^2}{2m}\int_{\mathbb R^d}|\nabla\psi|^2\,\dd^dq
  +
  \int_{\mathbb R^d}V(q)|\psi(q)|^2\,\dd^dq .
\]
For each \(x\in\mathbb R^d\) and \(\tau>0\), there is a countably additive
probability measure \(\mathbb W_x^\tau\) on
\[
  C([0,\tau],\mathbb R^d)
\]
whose finite-dimensional distributions are
\[
\begin{aligned}
  &\int F(\omega(t_1),\ldots,\omega(t_n))\,
  \dd\mathbb W_x^\tau(\omega)  \\
  &\quad =
  \int_{\mathbb R^{dn}}
  F(x_1,\ldots,x_n)
  \prod_{j=1}^n k^0_{t_j-t_{j-1}}(x_{j-1},x_j)\,
  \dd^dx_1\cdots\dd^dx_n ,
\end{aligned}
\]
for \(0<t_1<\cdots<t_n\le\tau\), with \(t_0=0\), \(x_0=x\), and
\[
  k^0_s(x,y)
  =
  \left(\frac{m}{2\pi\hbar s}\right)^{d/2}
  \exp\left[-\frac{m|x-y|^2}{2\hbar s}\right].
\]
This is Wiener measure with diffusion constant \(\hbar/m\).  For bounded
Borel \(f\), and then by extension for \(f\in L^2(\mathbb R^d)\),
\[
  \bigl(\ee^{-\tau H/\hbar}f\bigr)(x)
  =
  \int_{C([0,\tau],\mathbb R^d)}
  \exp\left[
    -\frac1\hbar\int_0^\tau V(\omega(s))\,\dd s
  \right]
  f(\omega(\tau))\,
  \dd\mathbb W_x^\tau(\omega)
\]
for almost every \(x\).  Disintegrating \(\mathbb W_x^\tau\) with respect to
the endpoint \(\omega(\tau)=y\) gives the Brownian-bridge probability measure
\(\mathbb W_{x,y}^\tau\), and the Euclidean transition kernel is
\[
  K_E(y,\tau;x,0)
  =
  k^0_\tau(x,y)
  \int
  \exp\left[
    -\frac1\hbar\int_0^\tau V(\omega(s))\,\dd s
  \right]
  \dd\mathbb W_{x,y}^\tau(\omega).
\]
\end{theorem}

\begin{proof}
The functions \(k^0_s(x,y)\) are nonnegative, integrate to one in \(y\), and
satisfy the convolution identity
\[
  \int_{\mathbb R^d} k^0_s(x,z)k^0_t(z,y)\,\dd^dz
  =
  k^0_{s+t}(x,y).
\]
The displayed finite-dimensional distributions are therefore consistent.
Kolmogorov extension gives a probability measure on paths with values in
\((\mathbb R^d)^{[0,\tau]}\).  The Gaussian increment moments obey
\[
  \mathbb E_{\mathbb W_x^\tau}
  |\omega(t)-\omega(s)|^{2r}
  =
  C_{d,r}\left(\frac{\hbar}{m}|t-s|\right)^r ,
\]
and Kolmogorov's continuity criterion gives a version supported on continuous
paths.  This proves the existence of the Wiener measure on
\(C([0,\tau],\mathbb R^d)\).

First assume that \(V\) is bounded and continuous.  Let
\(M_g\) denote multiplication by \(g\), and let
\[
  L=\frac{\hbar}{2m}\Delta .
\]
The heat semigroup \(\ee^{sL}\) has kernel \(k^0_s\).  Iterating the Markov
property gives
\[
\begin{aligned}
  &\bigl(\ee^{\delta L}M_{\exp(-\delta V/\hbar)}
  \bigr)^N f(x)\\
  &\quad =
  \mathbb E_x\left[
    \exp\left(
      -\frac{\delta}{\hbar}\sum_{j=1}^N V(B_{j\delta})
    \right)
    f(B_{N\delta})
  \right],
\end{aligned}
\]
where \(B\) is the coordinate process under \(\mathbb W_x^\tau\) and
\(\delta=\tau/N\).  Since Brownian paths are continuous and \(V\) is bounded
continuous, the Riemann sums converge almost surely to
\(\int_0^\tau V(B_s)\,\dd s\), and dominated convergence gives the Wiener
expectation.  The Trotter product formula identifies the same operator limit
with \(\ee^{-\tau H/\hbar}\), because
\[
  -\frac{H}{\hbar}=L-\frac1\hbar V
\]
as a form sum.

If \(V\) is continuous and bounded below, add a constant so that
\(V\ge0\), prove the formula for \(V_M=\min(V,M)\), and then let
\(M\to\infty\).  The path weights decrease pointwise to the displayed weight,
and the quadratic forms for \(T+V_M\) increase to the quadratic form for
\(T+V\); monotone convergence of closed forms gives strong convergence of the
semigroups.  The locally Kato bounded-below case is obtained by approximating
in the local Kato topology and using the Kato small-time estimate for the
negative part.  Finally, the bridge formula is the conditional-distribution
form of the preceding identity with respect to the endpoint density
\(k^0_\tau(x,y)\).
\end{proof}

Thus in ordinary Euclidean quantum mechanics the continuum path integral is a
genuine Wiener expectation on continuous paths, multiplied by the potential
weight, under the hypotheses of
Theorem~\ref{thm:wiener-feynman-kac-qm}.  The time-sliced Lebesgue products
are the approximating construction that converges to this measure-theoretic
object.
The Trotter product formula is still important, but its role is to identify
the operator semigroup represented by the Wiener integral; it is weaker than,
and should not be substituted for, the measure-level Feynman--Kac statement.
This theorem governs bosonic Schr\"odinger quantum mechanics, and by extension
supplies an important model for constructive scalar field theories when a
positive Euclidean measure exists.  The general path-integral framework is a
typed assignment of mathematical objects by a regulator and limiting
prescription.  Fermionic fields use Berezin integration over Grassmann
variables rather than countably additive Borel measures; gauge theories require
quotient, gauge-fixing, BRST/BV, or lattice data in addition to any measure on
fields; and topological terms such as a four-dimensional Yang--Mills
\(\theta\)-angle generally make the Euclidean weight complex.  Every use of
path-integral notation must therefore state the mathematical object being
used: a measure, an oscillatory integral, a Berezin functional, a
gauge-fixed/BV construction, a lattice regularization, a holonomy
construction, or a formal perturbative expansion.

\begin{theorem}[Faris--Lavine criterion]
\label{qthm:faris-lavine-essential-self-adjointness}
Essential self-adjointness on the algebraic core \(C_c^\infty(\mathbb R^d)\)
is a separate issue.  One convenient sufficient theorem is the Faris--Lavine
criterion, often cited as Reed--Simon Theorem X.41: if \(V\) is real,
\(C^\infty\), and satisfies
\[
  V(q)\ge -C(1+|q|^2)
\]
for some \(C<\infty\), then
\(-\Delta+V\) is essentially self-adjoint on \(C_c^\infty(\mathbb R^d)\).
Every smooth potential bounded below satisfies this growth-from-below
hypothesis, and polynomial oscillator potentials with positive leading part do
as well.  These cases have a unique self-adjoint closure from the natural
differential operator core, and the Trotter formula applies to that closure.
\end{theorem}

\begin{proof}
The criterion is a statement about the closure of a symmetric differential
operator determined on the algebraic core \(C_c^\infty(\mathbb R^d)\).
Write \(p_j=-\ii\partial_j\), \(p^2=-\Delta\), and first reduce by rescaling
and adding a constant to the model lower bound
\[
  V(x)\ge -|x|^2 .
\]
Let
\[
  A_0=p^2+V(x),
  \qquad
  N_0=p^2+V(x)+2|x|^2
\]
on \(C_c^\infty(\mathbb R^d)\).  Since \(V+2|x|^2\ge |x|^2\ge0\), the
Kato theorem for nonnegative locally square-integrable potentials gives a
positive self-adjoint closure \(N\) of \(N_0\).  The comparison with \(N\) is
the point of the proof.  On \(C_c^\infty(\mathbb R^d)\),
\[
  N_0-A_0=2|x|^2\ge0,
  \qquad
  N_0+A_0=2p^2+2(V+|x|^2)\ge0
\]
as quadratic forms, hence
\[
  -N_0\le A_0\le N_0
\]
in form order.  The oscillator estimate, obtained by integration by parts
from \(p^2+|x|^2\),
\[
  \||x|^2\psi\|
  \le
  C\bigl(\|N_0\psi\|+\|\psi\|\bigr),
  \qquad
  \psi\in C_c^\infty(\mathbb R^d),
\]
then implies that the \(N\)-graph norm controls \(A_0\).  Thus \(A_0\)
extends continuously from the \(N\)-graph domain, and one may apply Nelson's
commutator theorem with comparison operator \(N\).

It remains to check the commutator bound required by that theorem.  The only
noncommuting part is the added oscillator term:
\[
  \ii[N_0,A_0]
  =
  \ii[2|x|^2,p^2].
\]
Using \([x_j,p_k]=\ii\delta_{jk}\), one obtains, on the common algebraic
core,
\[
  \ii[N_0,A_0]
  =
  -4\sum_{j=1}^d(x_jp_j+p_jx_j).
\]
The form bound follows from completing the square:
\[
  (p_j\pm x_j)^*(p_j\pm x_j)
  =
  p_j^2+x_j^2\pm(p_jx_j+x_jp_j)\ge0,
\]
and therefore
\[
  \pm\sum_{j=1}^d(x_jp_j+p_jx_j)
  \le
  p^2+|x|^2
  \le
  N_0
\]
as quadratic forms.  Hence
\[
  \bigl|\langle\psi,\ii[N_0,A_0]\psi\rangle\bigr|
  \le
  4\,\langle\psi,N_0\psi\rangle .
\]
Nelson's theorem says that a symmetric operator \(A_0\) controlled in graph
norm by a positive self-adjoint \(N\), with the commutator bounded by
\(\langle\psi,N\psi\rangle\), is essentially self-adjoint on every
\(N\)-core contained in its domain.  Its proof uses the bounded regularized
operators
\[
  A_\epsilon=(1+\epsilon N)^{-1}A_0(1+\epsilon N)^{-1}
\]
and the commutator inequality to obtain a uniform Gronwall bound on the
\(N^{1/2}\)-energy of \(\ee^{\ii tA_\epsilon}\psi\); after passing
\(\epsilon\downarrow0\), the deficiency equations
\((A_0^*\mp\ii)\eta=0\) admit no nonzero \(L^2\) solutions.  This is the
operator-theoretic reason that the quadratic lower bound prevents a boundary
condition at infinity from being part of the Hamiltonian data.
\end{proof}

\begin{example}[A limit-circle failure of formal Hamiltonian data]
\label{ex:cubic-potential-limit-circle-failure}
The lower-growth hypothesis in the Faris--Lavine criterion is substantive.  In
one dimension, the symmetric differential operator
\[
  -\frac{\dd^2}{\dd q^2}+q^3
\]
on \(\mathcal S(\mathbb R)\) is a concrete failure mode.  At
\(q\to-\infty\), the deficiency equations
\[
  \left(-\frac{\dd^2}{\dd q^2}+q^3\right)\psi_\pm=\pm\ii\,\psi_\pm
\]
have two local solutions whose standard Liouville--Green asymptotics have
leading behavior
\[
  \psi_\pm(q)\sim |q|^{-3/4}
  \exp\left(\pm\ii\,\frac{2}{5}|q|^{5/2}\right).
\]
Both branches are square-integrable at that end because the squared amplitude
is asymptotic to \(|q|^{-3/2}\) and
\(\int^\infty r^{-3/2}\dd r<\infty\).  By the Weyl
limit-point/limit-circle classification theorem for one-dimensional
Sturm--Liouville operators, this places the endpoint \(q=-\infty\) in the
limit-circle case.  The symmetric operator is therefore not essentially
self-adjoint on the Schwartz core: a self-adjoint Hamiltonian requires an
additional boundary condition at this singular endpoint at infinity, and
different self-adjoint extensions give different unitary groups.  Thus the
formal expression
\(\widehat p^2/2m+V(\widehat q)\) does not by itself define the operator
appearing in \(\exp(-\ii T\widehat H/\hbar)\).
\end{example}

For the standard Hamiltonian
\[
  \widehat H=\frac{1}{2m}\widehat p^2+V(\widehat q),
\]
the product formula says that the operator obtained by alternating free
propagation and multiplication by the potential converges to the exact
evolution operator.  Inserting position resolutions between the factors gives
a finite-dimensional integral at each \(N\).  Thus the time-sliced path
integral is a representation of a product formula.

The distributional kernel
\[
  \bra{q_f}\exp\left(-\frac{\ii T}{\hbar}\widehat H\right)\ket{q_i}
\]
requires an additional topology.  The useful topologies include convergence
after smearing \(q_i\) and \(q_f\) with test functions, convergence after
analytic continuation to Euclidean time, and pointwise kernel convergence
under stronger regularity assumptions on \(V\).  Therefore the continuum
path-integral notation in this monograph always carries one of the following
meanings: a finite regulator, a strong operator limit, a smeared
distributional limit, a Euclidean measure or semigroup construction, or a
perturbative asymptotic expansion.

\section{Operator Ordering and Local Counterterms}

\begin{definition}[Regulated phase-space path-integral datum]
\label{def:regulated-phase-space-path-integral-datum}
The function \(h(q,p)\) depends on the ordering convention used to represent
\(\widehat H\). For example, using left endpoints \(q_n\) in the short-time
matrix element and using right endpoints \(q_{n+1}\) define different symbols
for the same operator. Their difference is of order \(\hbar\) in the
continuum action and appears as a local counterterm in the time-sliced
description.

At fixed \(N\), two representative choices are
\[
  \cdots
  \bra{q_{n+1}}p_n\rangle
  \bra{p_n}\ee^{-\ii\epsilon\widehat H/\hbar}\ket{q_n}
  \cdots,
  \qquad
  \bra{p_n}\widehat H\ket{q_n}
  =
  h(q_n,p_n)\braket{p_n}{q_n},
\]
and
\[
  \cdots
  \bra{p_{n+1}}q_{n+1}\rangle
  \bra{q_{n+1}}\ee^{-\ii\epsilon\widehat H/\hbar}\ket{p_n}
  \cdots,
  \qquad
  \bra{q_n}\widehat H\ket{p_n}
  =
  h'(q_n,p_n)\braket{q_n}{p_n}.
\]
They discretize the kinetic term respectively as
\[
  \sum_{n=0}^{N-1}p_{n,a}(q_{n+1}^a-q_n^a),
  \qquad
  \sum_{n=0}^{N-1}p_{n+1,a}(q_{n+1}^a-q_n^a),
\]
after relabelling endpoints where necessary.  The two symbols \(h\) and
\(h'\) represent the same operator only with their specified ordering
conventions, and they can differ by terms of order \(\hbar\).  Such terms
survive in the exponent as local regulator-dependent additions to the
effective Lagrangian.

Thus the data of the regulated path integral are
\[
  (\text{partition},\ \text{measure},\ \text{operator symbol},\
  \text{boundary conditions}).
\]
A formula involving \([Dq]\,[Dp]\) is meaningful as notation for a regulated
construction once these data have been declared.
\end{definition}

\section{Lagrangian Form for Quadratic Momentum Dependence}

\paragraph{Gaussian elimination of quadratic momenta.}
Assume the Hamiltonian symbol has the form
\[
  h(q,p)
  =
  \frac12\,G^{ab}(q)p_ap_b+V(q),
\]
where \(G^{ab}(q)\) is a positive-definite symmetric matrix and
\[
  G^{ab}(q)G_{bc}(q)=\delta^a_c.
\]
At fixed time slice, the momentum integral is Gaussian:
\[
  \int\frac{\dd^dp}{(2\pi\hbar)^d}
  \exp\left[
    \frac{\ii}{\hbar}
    \left(
      p_a\Delta q^a
      -\frac{\epsilon}{2}G^{ab}(q)p_ap_b
    \right)
  \right]
  =
  \frac{\sqrt{\det G_{ab}(q)}}{(2\pi\ii\hbar\epsilon)^{d/2}}
  \exp\left[
    \frac{\ii\epsilon}{\hbar}
    \frac12G_{ab}(q)v^av^b
  \right],
\]
where
\[
  \Delta q^a=q_{n+1}^a-q_n^a,
  \qquad
  v^a=\frac{\Delta q^a}{\epsilon}.
\]
The Lagrangian function obtained by Legendre transform is
\[
  L(q,\dot q)
  =
  \frac12G_{ab}(q)\dot q^a\dot q^b-V(q).
\]
Conversely, if a classical Lagrangian in generalized coordinates has this
metric form, as in the rigid-body configuration space example, then the
corresponding quantum Hamiltonian symbol is
\[
  h(q,p)=\frac12G^{ab}(q)p_ap_b+V(q)+O(\hbar)
\]
until an ordering convention and measure convention have been fixed.  The
\(O(\hbar)\) term records the same local regulator dependence described
above.
The finite-\(N\) configuration-space expression is therefore
\[
\begin{aligned}
  K_N(q_f,T;q_i,0)
  &=
  \int\prod_{n=1}^{N-1}\dd^dq_n
  \prod_{n=0}^{N-1}
  \frac{\sqrt{\det G_{ab}(q_n)}}{(2\pi\ii\hbar\epsilon)^{d/2}}
  \\
  &\quad\times
  \exp\left[
    \frac{\ii}{\hbar}
    \sum_{n=0}^{N-1}
    \epsilon\,
    L\left(q_n,\frac{q_{n+1}-q_n}{\epsilon}\right)
  \right].
\end{aligned}
\]
For \(G_{ab}=m\delta_{ab}\), the Lagrangian in this finite-dimensional
time-sliced kernel is
\[
  L(q,\dot q)=\frac{m}{2}\delta_{ab}\dot q^a\dot q^b-V(q).
\]

The calculation is the following finite-dimensional Gaussian completion.  At
fixed \(q_n\), one rewrites the exponent as
\[
  p_a\Delta q^a
  -\frac{\epsilon}{2}G^{ab}p_ap_b
  =
  -\frac{\epsilon}{2}G^{ab}
  \left(p_a-\frac1\epsilon G_{ac}\Delta q^c\right)
  \left(p_b-\frac1\epsilon G_{bd}\Delta q^d\right)
  +
  \frac{1}{2\epsilon}G_{ab}\Delta q^a\Delta q^b .
\]
The oscillatory Gaussian determinant gives
\(\sqrt{\det G_{ab}}(2\pi\ii\hbar\epsilon)^{-d/2}\), and the remaining
exponent is
\(\epsilon G_{ab}v^av^b/2\).  Multiplying over slices and including
\(-\epsilon V(q_n)\) gives the displayed configuration-space kernel.

\section{Time-Ordered Insertions and Sources}

\begin{definition}[Finite-slice time-ordered insertions and source kernel]
\label{def:finite-slice-time-ordered-insertions-source}
Let \(F(q(t_1),\ldots,q(t_r))\) be a product of coordinate insertions with
\[
  0<t_1<\cdots<t_r<T.
\]
At finite \(N\), choose the time slices nearest to the insertion times and
insert the corresponding coordinate functions into the finite-dimensional
integral. The resulting regulated expression represents
\[
  \bra{q_f}T\{\widehat q^{a_1}(t_1)\cdots
  \widehat q^{a_r}(t_r)\}U(T)\ket{q_i},
\]
where
\[
  \widehat q^a(t)
  =
  U(t)^{-1}\widehat q^aU(t)
\]
is the Heisenberg operator. The time ordering is inherited from the ordered
product of short-time evolution operators.

For a smooth compactly supported source \(J_a(t)\), define the finite-\(N\)
source factor by
\[
  \exp\left[
    \frac{\ii}{\hbar}
    \sum_{n=0}^{N-1}\epsilon\,J_a(t_n)q_n^a
  \right].
\]
The source-dependent kernel \(K_N[J]\) generates regulated time-ordered
coordinate insertions by differentiation with respect to \(J\). This is the
finite-dimensional ancestor of QFT generating functionals for Green
functions.
\end{definition}

\section{Hermitian Matrix Quantum Mechanics}
\label{sec:hermitian-matrix-quantum-mechanics}

Matrix quantum mechanics is a genuine \(0+1\)-dimensional quantum field
theory: the field at time \(t\) is a matrix, the symmetry may be global or
gauged, and the path integral is still a Hamiltonian time-evolution
construction of the kind developed above.  Its importance is not exhausted by
string interpretations.  The same system gives a controlled laboratory for
gauge projection, nontrivial configuration-space measure, collective
variables, and large-\(N\) limits.

\begin{definition}[Hermitian one-matrix quantum mechanics]
\label{def:hermitian-one-matrix-qm}
Let
\[
  \operatorname{Herm}_N=\{X\in M_N(\mathbb C):X^\dagger=X\}
\]
with flat metric
\[
  \|\delta X\|^2=\operatorname{tr}(\delta X)^2 .
\]
The kinematic Hilbert space is
\[
  \Hilb_N=L^2(\operatorname{Herm}_N,\dd X),
\]
where \(\dd X\) is the Lebesgue measure induced by this metric.  The global
\(U(N)\) action is
\begin{equation}
  X\longmapsto UXU^{-1},
  \qquad
  \Psi(X)\longmapsto \Psi(U^{-1}XU).
  \label{eq:matrix-qm-global-un-action}
\end{equation}
On the algebraic core \(C_c^\infty(\operatorname{Herm}_N)\), momenta are
\begin{equation}
  P_{ab}=-\ii\hbar\,\frac{\partial}{\partial X_{ba}},
  \qquad
  [X_{ab},P_{cd}]=\ii\hbar\,\delta_{ad}\delta_{bc},
  \label{eq:matrix-qm-canonical-commutator}
\end{equation}
with the Hermiticity relations understood in real coordinates on
\(\operatorname{Herm}_N\).  For a real potential \(V\), the formal invariant Hamiltonian is
\begin{equation}
  H_N
  =
  \operatorname{tr}
  \left(\frac12 P^2+V(X)\right)
  =
  -\frac{\hbar^2}{2}\Delta_{\operatorname{Herm}_N}
  +\operatorname{tr}V(X).
  \label{eq:hermitian-matrix-qm-hamiltonian}
\end{equation}
A self-adjoint Hamiltonian is the chosen closure or extension of
\eqref{eq:hermitian-matrix-qm-hamiltonian}; for confining polynomial
potentials with positive leading part, the Faris--Lavine mechanism above
gives the standard closure from the compactly supported smooth core.  Inverted
or unbounded potentials require an explicit self-adjointness and scattering
datum rather than a formal symbol alone.
\end{definition}

The gauged version is not obtained by first quantizing a string.  It is the
Hamiltonian system obtained by coupling the matrix coordinate to a
\(0+1\)-dimensional gauge field \(A_0(t)\) and then imposing the Gauss law.
For the single Hermitian matrix this amounts, at the Hilbert-space level, to
the invariant subspace
\[
  \Hilb_N^{\rm sing}
  =
  L^2(\operatorname{Herm}_N,\dd X)^{U(N)} .
\]
Equivalently, if
\[
  (G_\xi\Psi)(X)
  =
  -\ii\hbar\,\frac{\dd}{\dd s}\bigg|_{s=0}
  \Psi(\ee^{-s\xi}X\ee^{s\xi}),
  \qquad \xi\in\mathfrak u(N),
\]
then physical singlet states obey \(G_\xi\Psi=0\) for all \(\xi\).  This is
the finite-dimensional version of the gauge constraint: it removes gauge
orbits, not energy eigenstates by hand.

\paragraph{Eigenvalue coordinates and the collision set.}
On the regular locus
\[
  \operatorname{Herm}_N^{\rm reg}
  =
  \{X:\lambda_i(X)\ne\lambda_j(X)\ \text{for}\ i\ne j\},
\]
write
\[
  X=\Omega^{-1}\Lambda\Omega,\qquad
  \Lambda=\operatorname{diag}(\lambda_1,\ldots,\lambda_N).
\]
This description has a redundancy
\[
  (\Lambda,\Omega)
  \sim
  (W^{-1}\Lambda W,W^{-1}\Omega),
  \qquad W\in S_N,
\]
and a \(U(1)^N\) stabilizer of \(\Lambda\).  The metric identity
\begin{equation}
  \operatorname{tr}(\dd X)^2
  =
  \sum_{i=1}^N(\dd\lambda_i)^2
  +
  2\sum_{i<j}
  (\lambda_i-\lambda_j)^2
  |\theta_{ij}|^2,
  \qquad
  \theta=\dd\Omega\,\Omega^{-1},
  \label{eq:matrix-qm-eigenvalue-metric}
\end{equation}
gives the Jacobian
\begin{equation}
  \dd X
  =
  C_N\,\Delta(\lambda)^2
  \prod_{i=1}^N\dd\lambda_i\,\dd\mu(\Omega/T),
  \qquad
  \Delta(\lambda)=\prod_{i<j}(\lambda_i-\lambda_j),
  \label{eq:matrix-qm-vandermonde-measure}
\end{equation}
up to the conventional normalization of Haar measure on \(U(N)/T\).
The collision hyperplanes \(\lambda_i=\lambda_j\) are not removable
decorations.  They are the boundary of a Weyl chamber after quotienting by
\(S_N\), and domain conditions at this boundary are part of the radial
Hamiltonian.

For singlet wavefunctions
\(\Psi(X)=\widehat\Psi(\lambda_1,\ldots,\lambda_N)\), the flat Laplacian
reduces on the regular locus to
\begin{equation}
  \Delta_{\rm rad}\widehat\Psi
  =
  \frac1{\Delta^2}
  \sum_{i=1}^N
  \partial_i\!\left(\Delta^2\,\partial_i\widehat\Psi\right)
  =
  \sum_i\partial_i^2\widehat\Psi
  +
  2\sum_{i<j}
  \frac{\partial_i\widehat\Psi-\partial_j\widehat\Psi}
       {\lambda_i-\lambda_j}.
  \label{eq:matrix-qm-radial-laplacian}
\end{equation}
The singlet Hamiltonian is therefore symmetric on
\[
  L^2_{\rm sym}(\mathbb R^N,\Delta^2\prod_i\dd\lambda_i)
\]
with the potential \(\sum_i V(\lambda_i)\).  Multiplication by the
Vandermonde determinant is the crucial unitary change of variables:
\begin{equation}
  \psi(\lambda)=\Delta(\lambda)\widehat\Psi(\lambda).
  \label{eq:matrix-qm-vandermonde-wavefunction-map}
\end{equation}
Since \(\widehat\Psi\) is symmetric, \(\psi\) is antisymmetric.  Moreover,
away from the collision set,
\begin{equation}
  \Delta\,
  \Delta_{\rm rad}\widehat\Psi
  =
  \sum_{i=1}^N\partial_i^2(\Delta\widehat\Psi).
  \label{eq:matrix-qm-vandermonde-free-fermion-identity}
\end{equation}
Thus the gauged singlet sector is unitarily represented by \(N\) identical
one-dimensional fermions:
\begin{equation}
  H_{\rm sing}'
  =
  \sum_{i=1}^N
  \left(
    -\frac{\hbar^2}{2}\partial_{\lambda_i}^2+V(\lambda_i)
  \right)
  \quad
  \text{on antisymmetric wavefunctions.}
  \label{eq:matrix-qm-singlet-free-fermion-hamiltonian}
\end{equation}
Equivalently, on a single ordered Weyl chamber
\(\lambda_1>\lambda_2>\cdots>\lambda_N\), the fermionic wavefunction obeys
Dirichlet boundary conditions at \(\lambda_i=\lambda_j\).  This is the
operator meaning of the familiar statement that the Vandermonde makes the
singlet eigenvalues into free fermions.

\paragraph{Non-singlet angular sectors.}
The full matrix Hilbert space also contains non-singlet sectors.  After
Peter--Weyl decomposition on the angular variables, the radial Hamiltonian has
the representation-dependent form
\begin{equation}
  H_R'
  =
  \sum_i
  \left(
    -\frac{\hbar^2}{2}\partial_{\lambda_i}^2+V(\lambda_i)
  \right)
  +
  \frac{\hbar^2}{2}
  \sum_{i\ne j}
  \frac{\mathcal R_{ij}\mathcal R_{ji}}
       {(\lambda_i-\lambda_j)^2},
  \label{eq:matrix-qm-nonsinglet-radial-hamiltonian}
\end{equation}
where \(\mathcal R_{ij}\) act in the chosen angular representation.  The
inverse-square terms are not optional interpretation: they are the centrifugal
energy of changing the angular wavefunction as two eigenvalues approach.  A
later string or long-string comparison may attach additional language to
these sectors, but the intrinsic quantum-mechanical statement is simply the
representation-dependent radial problem
\eqref{eq:matrix-qm-nonsinglet-radial-hamiltonian}.

\paragraph{Large-\(N\) density and collective field.}
In the singlet sector define the eigenvalue density
\begin{equation}
  \rho(x)=\sum_{i=1}^N\delta(x-\lambda_i).
  \label{eq:matrix-qm-density-definition}
\end{equation}
At finite \(N\) this is an operator-valued distribution on the fermionic
Hilbert space, so replacing it by a smooth field is a large-\(N\) and
long-wavelength approximation.  The semiclassical phase-space description
uses the two Fermi-surface branches \(p_+(x)\ge p_-(x)\):
\begin{equation}
  \rho(x)=\frac{p_+(x)-p_-(x)}{2\pi\hbar},
  \qquad
  v(x)=\frac{p_+(x)+p_-(x)}2 .
  \label{eq:matrix-qm-fermi-branches-density-velocity}
\end{equation}
The leading hydrodynamic Hamiltonian follows by integrating the
single-particle energy over the filled phase-space strip,
\begin{align}
  H_{\rm coll}
  &=
  \int \dd x\,
  \left[
    \int_{p_-(x)}^{p_+(x)}
    \frac{\dd p}{2\pi\hbar}
    \left(\frac{p^2}{2}+V(x)\right)
  \right]
  \notag\\
  &=
  \int\dd x\,
  \left[
    \frac12\rho v^2
    +\frac{\pi^2\hbar^2}{6}\rho^3
    +V(x)\rho
  \right],
  \label{eq:matrix-qm-collective-hamiltonian}
\end{align}
with Poisson bracket
\[
  \{\rho(x),v(y)\}_{\rm P}
  =
  \partial_x\delta(x-y)
\]
when \(v=\partial_x\pi_\rho\) for the canonical field
\(\{\rho(x),\pi_\rho(y)\}_{\rm P}=\delta(x-y)\).  A more exact finite-\(N\)
collective-field quantization includes the Jacobian of the change from
fermion coordinates to density variables, gradient corrections, and boundary
terms at the edge of the Fermi sea.  The hydrodynamic formula
\eqref{eq:matrix-qm-collective-hamiltonian} is the leading controlled
coordinate, not a replacement for the underlying fermionic Hilbert space.

\begin{example}[Inverted oscillator and the \(c=1\) comparison boundary]
\label{ex:matrix-qm-inverted-oscillator-c1-boundary}
For
\[
  V(x)=-\frac12x^2
\]
the one-particle Hamiltonian is not confining.  A useful semiclassical state
fills one side of the inverted oscillator up to Fermi level \(-\mu\):
\[
  \frac12p^2-\frac12x^2<-\mu,\qquad x>\sqrt{2\mu}.
\]
The ground Fermi surface is
\begin{equation}
  p_\pm(x)=\pm\sqrt{x^2-2\mu},
  \qquad
  \rho_0(x)=\frac1{\pi\hbar}\sqrt{x^2-2\mu}.
  \label{eq:matrix-qm-inverted-density}
\end{equation}
Writing \(x=\sqrt{2\mu}\cosh\tau\), the asymptotic region
\(\tau\to\infty\) turns the quadratic part of the collective Hamiltonian into
that of a massless scalar on a half-line, while the tip
\(\tau=0\) requires boundary and counterterm data.  This is already a
complete quantum-mechanical scattering problem for fermions in an inverted
potential.  The double-scaled \(c=1\) string interpretation, particle-hole
leg factors, nonperturbative ZZ-instanton language, and JT-gravity
matrix-model analogies are comparison statements attached to this quantum
mechanics; they are not used here as foundations for
\eqref{eq:matrix-qm-radial-laplacian}--
\eqref{eq:matrix-qm-collective-hamiltonian}.
\end{example}

\section{Euclidean Traces and Periodic Boundary Conditions}

\paragraph{Trace and periodic Euclidean boundary conditions.}
Assume that \(\widehat H\) is semibounded and that
\(\ee^{-\beta_{\mathrm T}\widehat H/\hbar}\) is trace class for
\(\beta_{\mathrm T}>0\). The Euclidean thermal trace is
\[
  Z(\beta_{\mathrm T})
  =
  \operatorname{Tr}_{\Hilb_{\mathcal Q}}
  \ee^{-\beta_{\mathrm T}\widehat H/\hbar}.
\]
This trace-class hypothesis is a finite-volume or confining
quantum-mechanical hypothesis.  It is the correct assumption for a harmonic
oscillator, for stable anharmonic oscillators with discrete spectrum growing
fast enough, and for field theories after both an ultraviolet regulator and a
finite spatial-volume regulator have been imposed.  It is not the Hilbert-space
structure of an infinite-volume translation-invariant QFT.  In infinite
spatial volume, translating a finite-energy local excitation through disjoint
spatial regions produces the usual volume divergence; correspondingly
\(\operatorname{Tr}\ee^{-\beta H}\) is infinite even when local thermal
correlators exist.

Thermal QFT is therefore formulated by one of two equivalent local procedures.
One may put the theory in a finite spatial box \(\Lambda\), form
\[
  \rho_{\beta,\Lambda}
  =
  Z_{\beta,\Lambda}^{-1}\ee^{-\beta H_\Lambda},
  \qquad
  Z_{\beta,\Lambda}=\operatorname{Tr}_{\Hilb_\Lambda}\ee^{-\beta H_\Lambda},
\]
and then take thermodynamic limits of local observables.  Or one may define
the infinite-volume thermal state directly as a KMS state on the quasilocal
observable algebra, in the sense of
Definition~\ref{def:kms-state}.  Formal traces per unit volume, or traces with
the translation volume divided out, compute densities rather than the literal
trace over the full infinite-volume Hilbert space.
For a trace-class quantum-mechanical Hamiltonian, the trace of the heat
operator is obtained by integrating its diagonal position kernel:
\[
  Z(\beta_{\mathrm T})
  =
  \int_{\mathbb R^d}\dd^dq_0\,
  \bra{q_0}
  \ee^{-\beta_{\mathrm T}\widehat H/\hbar}
  \ket{q_0}.
\]
The same time-sliced construction as above, now with Euclidean time
\(\tau\in[0,\beta_{\mathrm T}]\), therefore has
\[
  q_N=q_0 .
\]
Thus the trace of a bosonic quantum-mechanical system is represented by the
Euclidean path integral over periodic paths
\[
  q(\tau+\beta_{\mathrm T})=q(\tau).
\]
This conclusion is a statement about the boundary condition produced by the
trace.  When Theorem~\ref{thm:wiener-feynman-kac-qm} applies, the
corresponding continuum object is the Feynman--Kac weight integrated over
Brownian loops, obtained by disintegrating Wiener measure with the endpoint
constraint
\[
  q(\beta_{\rm T})=q(0).
\]
Outside such hypotheses, the trace path integral remains a regulated or formal
expression whose convergence must be specified separately.

\section{Vacuum Projection}

\paragraph{Euclidean long-time projection to a gapped ground state.}
Assume \(\widehat H\) has a normalized ground state \(\ket{0}\), with energy
\(E_0\), and that a vector \(\ket\psi\in\Hilb_{\mathcal Q}\) satisfies
\(\braket{0}{\psi}\ne0\). If the spectrum above \(E_0\) is separated by a
gap, then
\[
  \lim_{\tau\to\infty}
  \frac{\ee^{-\tau(\widehat H-E_0)/\hbar}\ket\psi}
       {\left\|\ee^{-\tau(\widehat H-E_0)/\hbar}\ket\psi\right\|}
  =
  \ket{0}
\]
up to the phase of \(\braket{0}{\psi}\). Vacuum correlation functions can
therefore be obtained from long Euclidean-time evolution:
\[
  \bra{0}T\{\widehat q^{a_1}(t_1)\cdots
  \widehat q^{a_r}(t_r)\}\ket{0}
\]
is the normalized long-time limit of the corresponding transition amplitude
with insertions, after the Euclidean projection has been applied.
This is the elementary spectral projection mechanism behind Euclidean vacuum
boundary conditions.  By the spectral theorem for \(\widehat H\), write
\(\ket\psi=c_0\ket0+\eta\) with \(c_0=\braket0\psi\ne0\) and \(\eta\)
orthogonal to the ground-state subspace.  The gap assumption gives
\[
  \|\ee^{-\tau(\widehat H-E_0)/\hbar}\eta\|
  \le
  \ee^{-\tau\Delta/\hbar}\|\eta\|
\]
for some \(\Delta>0\).  Therefore the normalized Euclidean evolution tends to
\(\ee^{\ii\arg c_0}\ket0\).  Applying the same spectral projection to initial
and final wave packets in a transition amplitude with time-ordered insertions
selects the vacuum matrix element after the common exponential normalization is
removed.
The substantive hypotheses are the existence of the normalized ground state,
the nonzero overlap of the boundary vector with it, and the gap.  Without a
gap the same formula can fail in norm, and one must specify the weaker
spectral or distributional limit being used.

The next chapter will analyze the analytic continuation and Gaussian
integrals that make this construction computationally useful. The present
chapter has introduced the regulated path-integral data and its relation to
Hamiltonian time evolution.
